% Secao 3 - Dados e construcao. Reaproveita convencoes do paper18.
% Disciplina: descritivo. Numeros com % src.

\section{Dados e construção}
\label{sec:dados}

O relatório combina quatro fontes para o período 2015--2022, todas no nível do
município de residência. Os fluxos de pacientes vêm do Sistema de Informações
Hospitalares (SIH/DATASUS): cada Autorização de Internação Hospitalar (AIH)
registra o município de residência do paciente e o estabelecimento (CNES) onde a
internação ocorreu, o que permite construir o grafo bipartido
residência~$\rightarrow$~hospital que define o fluxo. O universo de hospitais
ativos por ano vem do Cadastro Nacional de Estabelecimentos de Saúde (CNES),
restrito a estabelecimentos com leito hospitalar instalado. As distâncias e os
tempos de viagem usam os centroides municipais do IBGE e a malha rodoviária via
OSRM. A população, denominador das taxas, vem das estimativas IBGE/SIDRA.

Duas disciplinas de construção merecem registro. A primeira é o \emph{firewall}
entre residência e provedor: toda exposição é atribuída pelo município onde o
paciente \emph{mora}, nunca pelo município onde o hospital opera, de modo que o
fluxo mede a dependência revelada do morador, não a atividade do hospital. A
segunda são as armadilhas conhecidas do SIH: contamos AIH de internação (AIH-1)
e não as de continuação (AIH-5), para não duplicar episódios, e tratamos o SIH
como retrato do SUS --- internações na saúde suplementar não aparecem, o que
subestima a oferta efetiva em municípios de alta penetração privada, ressalva
que a Seção~\ref{sec:divergencia} incorpora via cobertura ANS.

A amostra cobre 5.588 municípios de residência, 10.450 hospitais com leito ativos
na janela e 94,0 milhões de internações.% src: bloco cobertura - n_mun_res 5588, n_hosp_beds 10450, n_aih 93.996.098
O recorte analítico do relatório usa os 5.492 municípios com fluxo suficiente
para um burden efetivo bem definido.% src: master_muni_panel N=5492 (apos filtro NaN travel_burden)
